\scanpage{11}\mspage{10}
En \deleted{\unsure{l}}effet, c'est connu pour les sous-familles admissibles d'une famille (ici $\unsure{\mathscr{L}|\Phi}\cup\deleted{\unsure{\mathscr{L}}}\Psi$) -- cf. Scholie \unsure{pge|pg} 7.

\margin{\uline{Cor} 3 Si $\mathscr{L},\Psi$ comp., alors pour $\mathscr{L}',\deleted{\unsure{\Psi}}\leq\mathscr{L},\ \Psi'\leq\Psi$ (parties fermées) $\mathscr{L}'$ et $\Psi'$, st admissibles, et $\mathscr{L}'\cup\Psi'$ est fermé (ss-famille adm.) dans $\mathscr{L}\cup\Psi$.}

% slanted
\margin{On appellera ainsi les éléments de $\unsure{\mathcal F|\mathscr F}$ les \uline{``figures ensemblistes''} de $L$}

Soit $\mathcal F$ l'ens des \deleted{toutes} les parties \uline{admissibles} $\mathscr{L}$ de $\unsure{\mathcal P|\mathfrak P}(L)$. On ordonne $\mathcal F$ par
\[
\begin{aligned}
\mathscr{L}\leq\Psi\quad\overset{\mathrm{def}}{\Longleftrightarrow}\quad &\mathscr{L}\ \text{est une ss-}\\
&\text{famille admissible}\\
&\text{de }\Psi.
\end{aligned}
\]

Alors on a

\uline{C 1} $\mathcal F$ admet des Sup majorés, car si $\mathscr{L}_i\in\mathcal F,\ \mathscr{L}_i\leq\deleted{\mathscr{L}}\Psi$ des ss-ens. fermés d'un $=$ $\deleted{\mathscr{L}}\Psi$, alors $\bigcup_i\mathscr{L}_i$ est un sous-ens. fermé de $\Psi$, qui est le la sup. C'est immédiat par \unsure{le sorite|les sorites} sur les sous-familles \unsure{équiv.|égaux.} (\unsure{de ses|i.e.|les|deux} parties fermées) d'un $\Psi$

% slanted
\margin{La Sup pour la relation $\leq$ (qui est ici $\subset$) est $\bigcup_i$}

\deleted{\unsure{Oui|Or|O.K.}, un peu \unsure{vif|cf.}, il faut \unsure{passer|préciser} par des Sup finis, en notant que $\emptyset\subset\mathcal P(L)$ est admissible (! cas du Sup \inserted{de la fam.} vide !) -- la seule famille admissible $\underset{\in\mathcal F}{\mathscr{L}}$ telle que $\lvert\mathscr{L}\rvert=\emptyset$ et que si $\mathscr{L},\Psi,\Sigma$ sont tels que}

% slanted
\margin{\unsure{explicitons|expliciter} \unsure{cf.} \uuline{Prop} ou Cor}

% This line has a horizontal cancellation in addition to the diagonal strokes across the draft.
\deleted{$\mathscr{L}$ et $\Psi$) est admissible, alors $\Sigma$ \unsure{ai|est}}

\deleted{($\mathscr{L},\Psi$) compatibles, alors $\Sigma$ est compatible avec $\mathscr{L}\unsure{\cup|et}\Psi$ \unsure{sss|ssi} il l'est avec $\mathscr{L}$ et avec $\Psi$}
